\section{Ce que les trente-cinq trades racontent}

\lettrine{L}{e} cycle d'expansion momentum 2026-H2 cherchait des instruments
capables de porter le même moteur que la sleeve or, pour cesser de faire
dépendre 80\,\% du résultat du portefeuille de trente-cinq trades sur un seul
sous-jacent. Vingt-et-un instruments ont été passés au simulateur~; USD/JPY en
est ressorti premier. Ce document regarde ses transactions une par une, parce
qu'un classement par ratio ne dit pas \emph{d'où} vient le résultat --- et sur
ce candidat, la réponse est inattendue.

Le matériau est le journal d'exécution d'un \textbf{run de recherche} MT5~: la
sleeve TSMOM jouée seule, à pleine allocation, facteur de risque~1,0, sur un
dépôt de 100\,000~dollars, du 1\textsuperscript{er}~janvier~2021 au
30~avril~2026. Trente-cinq positions y figurent, du 14~septembre~2021 au
29~avril~2026, avec pour chacune le prix d'entrée, le prix de sortie, le volume
en lots, le résultat, le portage, ainsi que le score de décision et le levier
relevés dans le journal du simulateur. Tous les chiffres de ce document
proviennent de cette source unique, et le total recolle au cent près sur le
rapport du simulateur~: $+$176\,822,48\,\$.

\subsection*{Cinq constats}

\begin{enumerate}[leftmargin=1.4em, itemsep=0.45em]

\item \textbf{Le portage fait les deux tiers du résultat.} Le mouvement du
change n'a rapporté que $+$64\,557\,\$~; le différentiel de taux encaissé sur
les positions en a ajouté $+$112\,266\,\$, soit \metric{63\,\%} du net. C'est
l'exact miroir de la sleeve or, où le portage \emph{coûte} 28\,\% du résultat
brut. Ici, être long dollar contre yen, c'est être payé pour attendre.

\item \textbf{Une seule position fait tout --- et davantage.} Le trade ouvert
le 24~septembre~2021 et tenu \metric{416~jours} rapporte $+$188\,374\,\$, soit
\emph{plus} que le résultat total. Retiré, il ne reste pas un gain plus modeste
mais une perte de $-$11\,551\,\$ sur les trente-quatre autres. La concentration
est ici plus extrême que sur l'or.

\item \textbf{Le taux de réussite est trompeusement équilibré.} Dix-huit trades
sur trente-cinq finissent positifs, soit 51,4\,\%. Mais seulement
\metric{quinze} sont gagnants \emph{sur le prix}~: le portage retourne le signe
de trois d'entre eux. Le gain moyen ($+$25\,744\,\$) ne vaut que
\metric{1,5~fois} la perte moyenne ($-$16\,857\,\$) --- l'asymétrie qui fait
vivre l'or est ici bien plus faible.

\item \textbf{Les mauvaises passes sont courtes mais violentes.} La plus longue
série perdante ne compte que \textbf{trois} trades, contre dix sur l'or. Mais
ces trois clôtures-là, du 21~mars au 3~mai~2023, coûtent $-$96\,940\,\$ ---
c'est exactement le repli maximal de balance mesuré par le simulateur,
\metric{33,8\,\%} du compte, encaissé en six semaines.

\item \textbf{Le plafond de levier est la contrainte réelle.} La volatilité du
yen évolue entre 2,9 et 18,1\,\% annualisés sur la période, médiane 8,2\,\%,
très loin de la cible de 55\,\%. Le levier sature donc à 6,6$\times$ dès que la
volatilité passe sous 8,3\,\% --- ce qui arrive une séance sur deux, et pour
\metric{18~entrées sur 35}.

\end{enumerate}

\begin{warningbox}
\textbf{Ce document décrit un candidat, pas une sleeve déployée.} Le run est un
run de recherche~: sleeve isolée, allocation 100\,\%, facteur de risque~1,0,
dépôt de 100\,000\,\$. Les montants ne sont comparables ni à ceux du
portefeuille client ni à ceux de la sleeve or, qui tourne à 10\,\% d'allocation
avec un facteur de risque de 4,5. Le simulateur a par ailleurs tourné en mode
\emph{OHLC~1~minute}, qui interpole les exécutions et flatte les résultats.
Trente-cinq observations ne permettent aucune conclusion statistique par
sous-groupe. Le chapitre~\ref{sec:limites} détaille ce que ce document ne
démontre pas.
\end{warningbox}

% Généré par scripts/build_usdjpy_trades_figures.py — NE PAS ÉDITER À LA MAIN.
\begin{table}[H]
\centering
\begin{tabular}{@{}lr@{}}
\toprule
\textbf{Métrique} & \textbf{Run MT5 de recherche} \\
\midrule
Trades & 35 \\
Gagnants & 18 (51.4\,\%) \\
Gagnants avant portage & 15 \\
Résultat net & +176,822\,\$ \\
Résultat de prix & +64,557\,\$ \\
Portage encaissé & +112,266\,\$ \\
Gain moyen & +25,744\,\$ \\
Perte moyenne & -16,857\,\$ \\
Meilleur trade & +188,374\,\$ \\
Pire trade & -58,262\,\$ \\
Résultat médian & +273\,\$ \\
Rendement de prix moyen & +0.51\,\% \\
Durée médiane & 5\,j \\
Durée maximale & 416\,j \\
Repli maximal de balance & 33.8\,\% \\
Stops de sécurité & 0 \\
\bottomrule
\end{tabular}
\caption{Le candidat USD/JPY en chiffres, sleeve isolée sur un dépôt de 100\,000\,\$, période 2021--2026.}
\end{table}

